% ============================================================
% esp32_pictorial.tex — spatial / colour-coded build view.
%
% Complementary to esp32_supervisor.svg, not a substitute. This
% sheet answers "which coloured wire goes where on the bench",
% laid out the way a breadboard actually sits: two power rails
% along the bottom, modules above them.
%
% Pin-level truth lives in WIRING.md. Signal topology lives in
% esp32_supervisor.svg. This drawing is spatial only.
%
% Routing rule: vertical wires pass OVER horizontal ones (white
% under-stroke). A filled dot is the only thing that ever means
% "connected".
%
% Regenerate with:  make esp32_pictorial.svg
% ============================================================
\def\pgfsysdriver{pgfsys-dvisvgm.def}
\documentclass[border=10pt,dvisvgm]{standalone}
% ============================================================
% tsstyle.tex — shared preamble for every sheet in this directory.
%
% Inputted by each sheet after \documentclass. Sheets stay pure
% geometry; everything about how the geometry *looks* lives here,
% so a style change is a one-file edit.
% ============================================================

\usepackage[T1]{fontenc}
\usepackage[utf8]{inputenc}
\usepackage[scaled=0.92]{helvet}
\renewcommand{\familydefault}{\sfdefault}

\usepackage[american,siunitx,RPvoltages]{circuitikz}
\usetikzlibrary{calc,positioning,fit,backgrounds,arrows.meta,decorations.pathreplacing}

% ------------------------------------------------------------
% Palette. Deliberately small: ink for conductors, mute for
% annotation, and exactly two semantic colours — mains (danger)
% and SELV (safe). Colour carries meaning on these sheets; it is
% never decoration.
% ------------------------------------------------------------
\definecolor{tsink}{HTML}{16181D}    % conductors, symbol bodies
\definecolor{tsmute}{HTML}{6E7178}   % secondary annotation
\definecolor{tsrule}{HTML}{C7CAD1}   % boxes, leader lines
\definecolor{tsmains}{HTML}{B3261E}  % mains potential — danger
\definecolor{tsselv}{HTML}{15607A}   % SELV / low-voltage rail
\definecolor{tssafe}{HTML}{1E6B45}   % fail-safe / passive layer
\definecolor{tspanel}{HTML}{F5F6F7}  % block fills
\definecolor{tspanelb}{HTML}{E4E7EB} % block fills, alternate

% ------------------------------------------------------------
% Global circuitikz sizing. Bipoles are scaled down from the
% default because these sheets carry a lot of elements per unit
% area; the default American resistor is comically large next to
% an IC block.
% ------------------------------------------------------------
\ctikzset{
  bipoles/length=0.85cm,
  resistors/scale=0.85,
  capacitors/scale=0.8,
  label/align=straight,
  font=\footnotesize,
}

\tikzset{
  % Conductors -------------------------------------------------
  wire/.style        = {tsink, line width=0.9pt, line cap=round},
  buswire/.style     = {tsink, line width=1.5pt, line cap=round},
  mainswire/.style   = {tsmains, line width=1.5pt, line cap=round},
  selvwire/.style    = {tsselv, line width=1.2pt, line cap=round},
  % Functional blocks -----------------------------------------
  block/.style       = {draw=tsink, line width=1pt, fill=tspanel,
                        rounded corners=1.5pt, align=center,
                        inner sep=6pt, font=\small\bfseries},
  subblock/.style    = {draw=tsrule, line width=0.7pt, fill=white,
                        rounded corners=1.5pt, align=center,
                        inner sep=4pt, font=\scriptsize},
  % Text ------------------------------------------------------
  pinlbl/.style      = {font=\scriptsize, tsink, inner sep=2pt},
  siglbl/.style      = {font=\scriptsize\itshape, tsmute, inner sep=1.5pt},
  netlbl/.style      = {font=\scriptsize\bfseries, tsink, inner sep=2pt},
  note/.style        = {font=\scriptsize, tsmute, align=left},
  danger/.style      = {font=\scriptsize\bfseries, tsmains, align=left},
  % Isolation boundary ----------------------------------------
  isoline/.style     = {tsmains, line width=1pt, dash pattern=on 5pt off 3pt},
  % A wire that crosses another without connecting. The white
  % under-stroke reads as an over-pass, so a crossing can never be
  % mistaken for a junction (junctions are always a filled dot).
  overpass/.style    = {preaction={draw, white, line width=3.4pt}},
  % Sheet furniture -------------------------------------------
  titleblock/.style  = {draw=tsink, line width=1pt, fill=white,
                        align=left, inner sep=7pt, font=\scriptsize},
  legendbox/.style   = {draw=tsrule, line width=0.7pt, fill=white,
                        align=left, inner sep=7pt, font=\scriptsize},
}

% ------------------------------------------------------------
% \pinrow{<x>}{<y>}{<dir>}{<label>}  — a pin stub on a block
% edge. dir is -1 (left edge, stub runs left) or +1 (right edge).
% ------------------------------------------------------------
\newcommand{\pinstub}[5]{%
  \draw[wire] (#1,#2) -- ++(#3*0.55,0);
  \node[pinlbl, anchor=#4] at (#1+#3*0.10,#2+0.14) {#5};
}

% ------------------------------------------------------------
% \titleblocknode{<x>}{<y>}{<sheet>}{<subtitle>}{<source>}
% Standard title block, anchored at its south east corner.
% ------------------------------------------------------------
\newcommand{\sheettitle}[5]{%
  \node[titleblock, anchor=south east] at (#1,#2) {%
    {\normalsize\bfseries #3}\\[2pt]
    {\color{tsmute}#4}\\[4pt]
    {\color{tsmute}Table-saw thermal protection retrofit\hfill}\\
    {\color{tsmute}Generated from \texttt{#5} — do not edit the SVG}%
  };
}


% ------------------------------------------------------------
% Breadboard wire colour code — matches WIRING.md exactly.
% Yellow and grey are darkened from the physical wire colour so
% they stay legible as 1pt strokes on white.
% ------------------------------------------------------------
\definecolor{wred}{HTML}{CE3B2F}
\definecolor{worange}{HTML}{D9741A}
\definecolor{wblack}{HTML}{2B2B2B}
\definecolor{wyellow}{HTML}{B08900}
\definecolor{wgreen}{HTML}{2C8A44}
\definecolor{wblue}{HTML}{1E6FB8}
\definecolor{wpurple}{HTML}{7A4CA8}
\definecolor{wbrown}{HTML}{8A5426}
\definecolor{wgrey}{HTML}{75797F}

\tikzset{
  w/.style   = {line width=1.3pt, line cap=round, line join=round},
  vw/.style  = {w, overpass},   % vertical run — passes over horizontals
}

\begin{document}
\begin{circuitikz}[line width=0.9pt]

% ============================================================
% Provenance banner — which hardware this sheet is drawn for
% ============================================================
\node[font=\scriptsize\bfseries, tsmains, anchor=south west] at (-10.9,10.6)
  {PATH B TARGET --- MAX31855, RELAY MODULE AND ISOLATED PSU ARE
   \textbf{NOT PURCHASED} (TASK-1/2/3).};
\node[font=\scriptsize, tsmute, anchor=south west] at (-10.9,10.2)
  {For the hardware actually on hand, build from
   \texttt{pathA\_supervisor.svg} and \texttt{WIRING-PATH-A.md}.};

% ============================================================
% Power rails (breadboard convention: rails along the bottom)
% ============================================================
\draw[w, wred]   (-5.4,-1.6) -- (9.6,-1.6);
\draw[w, wblack] (-5.4,-2.4) -- (9.6,-2.4);
\node[font=\scriptsize\bfseries, wred,   anchor=west] at (9.75,-1.6) {+5\,V rail};
\node[font=\scriptsize\bfseries, wblack, anchor=west] at (9.75,-2.4) {GND rail};

% ============================================================
% Isolated AC-DC PSU
% ============================================================
\draw[block] (-9.0,-3.4) rectangle (-5.4,-0.6);
\node[font=\small\bfseries, anchor=north, align=center] at (-7.2,-1.05)
  {Isolated\\AC-DC PSU};
\node[note, anchor=north, align=center] at (-7.2,-2.35)
  {Mean Well\\IRM-05-5};
\node[pinlbl, anchor=east] at (-5.52,-1.6) {+5V};
\node[pinlbl, anchor=east] at (-5.52,-2.4) {$-$V};
\draw[w, wred]   (-5.4,-1.6) -- (-5.4,-1.6);
\draw (-5.4,-1.6) node[circ, color=wred]{};
\draw (-5.4,-2.4) node[circ, color=wblack]{};

% mains input — the only mains-potential wiring on this sheet
\draw[mainswire] (-9.0,-1.6) -- (-10.6,-1.6);
\draw[mainswire] (-9.0,-2.4) -- (-10.6,-2.4);
\node[pinlbl, anchor=west] at (-8.88,-1.6) {L};
\node[pinlbl, anchor=west] at (-8.88,-2.4) {N};
\node[danger, anchor=east, align=right] at (-10.75,-2.0)
  {MAINS\\from the\\250\,mA fuse};

% ============================================================
% ACK button
% ============================================================
\draw[block, fill=white] (-9.0,2.0) rectangle (-5.4,3.8);
\node[font=\small\bfseries, anchor=north] at (-7.2,3.75) {ACK button};
\node[note, anchor=north, align=center] at (-7.2,3.15)
  {momentary NO\\panel mount};
\node[pinlbl, anchor=east] at (-5.52,3.0) {pole 1};
\node[pinlbl, anchor=east] at (-5.52,2.4) {pole 2};

\draw[w, wgrey]  (-5.4,3.0) -- (0,3.0);
\draw[w, wblack] (-5.4,2.4) -- (-5.0,2.4);
\draw[vw, wblack] (-5.0,2.4) -- (-5.0,-2.4);
\draw (-5.0,-2.4) node[circ, color=wblack]{};

% ============================================================
% ESP32-DevKitC-32E
% ============================================================
\draw[block, fill=tspanelb] (0,1.6) rectangle (5,9.8);
\node[font=\small\bfseries, anchor=north, align=center] at (2.5,7.1)
  {ESP32-\\DevKitC-32E};
\node[note, anchor=north, align=center] at (2.5,5.6)
  {\texttt{B0GF1ZJCCN}\\38-pin, USB-C\\3.3\,V logic};
\node[note, anchor=north, align=center, tssafe] at (2.5,3.5)
  {GPIO2 drives the\\onboard blue LED —\\no external wire};

% left edge
\node[pinlbl, anchor=west] at (0.12,8.6) {GND};
\node[pinlbl, anchor=west] at (0.12,7.8) {VIN};
\node[pinlbl, anchor=west] at (0.12,3.0) {GPIO27};
% right edge
\node[pinlbl, anchor=east] at (4.88,9.0) {GPIO18};
\node[pinlbl, anchor=east] at (4.88,8.4) {GPIO19};
\node[pinlbl, anchor=east] at (4.88,7.8) {GPIO5};
\node[pinlbl, anchor=east] at (4.88,6.4) {3V3};
\node[pinlbl, anchor=east] at (4.88,4.6) {GPIO34};
\node[pinlbl, anchor=east] at (4.88,3.6) {GPIO26};

% ---- ESP32 power down to the rails ----
\draw[w, wblack]  (0,8.6) -- (-2.0,8.6);
\draw[vw, wblack] (-2.0,8.6) -- (-2.0,-2.4);
\draw (-2.0,-2.4) node[circ, color=wblack]{};
\draw[w, wred]    (0,7.8) -- (-1.2,7.8);
\draw[vw, wred]   (-1.2,7.8) -- (-1.2,-1.6);
\draw (-1.2,-1.6) node[circ, color=wred]{};

% ============================================================
% SPI bus — three straight runs to the MAX31855
% ============================================================
\draw[w, wyellow] (5,9.0) -- (9.5,9.0);
\draw[w, wgreen]  (5,8.4) -- (9.5,8.4);
\draw[w, wblue]   (5,7.8) -- (9.5,7.8);

% ============================================================
% MAX31855 breakout
% ============================================================
\draw[block, fill=white] (9.5,6.0) rectangle (13.5,9.8);
\node[font=\small\bfseries, anchor=north] at (11.5,9.7) {MAX31855};
\node[note, anchor=north, align=center] at (11.5,6.9)
  {K-type breakout\\3.3\,V only};
\node[pinlbl, anchor=west] at (9.62,9.0) {SCK};
\node[pinlbl, anchor=west] at (9.62,8.4) {SO};
\node[pinlbl, anchor=west] at (9.62,7.8) {CS};
\node[pinlbl, anchor=west] at (9.62,7.2) {VCC};
\node[pinlbl, anchor=west] at (9.62,6.6) {GND};
\node[pinlbl, anchor=east] at (13.38,8.8) {T+};
\node[pinlbl, anchor=east] at (13.38,8.0) {T$-$};

% thermocouple out to the motor pigtail
\draw[w, wyellow] (13.5,8.8) -- (15.4,8.8);
\draw[w, wred]    (13.5,8.0) -- (15.4,8.0);
\draw (15.4,8.8) node[ocirc, color=wyellow]{};
\draw (15.4,8.0) node[ocirc, color=wred]{};
\node[note, anchor=west] at (15.6,8.8) {K-type \textbf{yellow} $(+)$};
\node[note, anchor=west] at (15.6,8.0) {K-type \textbf{red} $(-)$};
\node[note, anchor=west, align=left, tssafe] at (15.6,7.15)
  {to the winding, via\\\texttt{../harness/motor\_pigtail.yml}};

% MAX31855 GND down to the rail
\draw[w, wblack]  (9.5,6.6) -- (9.1,6.6);
\draw[vw, wblack] (9.1,6.6) -- (9.1,-2.4);
\draw (9.1,-2.4) node[circ, color=wblack]{};

% ============================================================
% 3V3 rail — ESP32 regulator out to the NTC divider + MAX31855
% ============================================================
\draw[w, worange]  (5,6.4) -- (8.9,6.4);
\draw[vw, worange] (8.9,6.4) -- (8.9,7.2);
\draw[w, worange]  (8.9,7.2) -- (9.5,7.2);
\draw[vw, worange] (7.5,6.4) -- (7.5,6.0);
\draw (7.5,6.4) node[circ, color=worange]{};

% ============================================================
% NTC frame divider
% ============================================================
\draw[subblock] (6.4,4.0) rectangle (8.6,6.0);
\node[font=\scriptsize\bfseries, anchor=north] at (7.5,5.85) {NTC divider};
\node[note, anchor=north, align=center] at (7.5,5.4)
  {10\,k$\Omega$ 1\% pullup\\$+$ NTC 10\,k$\Omega$\\B3950 on frame};
\draw[w, wpurple]  (5,4.6) -- (6.4,4.6);
\draw[vw, wblack]  (7.5,4.0) -- (7.5,-2.4);
\draw (7.5,-2.4) node[circ, color=wblack]{};

% ============================================================
% Relay module
% ============================================================
\draw[w, wbrown]  (5,3.6) -- (9.3,3.6);
\draw[vw, wbrown] (9.3,3.6) -- (9.3,3.0);
\draw[w, wbrown]  (9.3,3.0) -- (9.5,3.0);

% Header order is IN / GND / VCC rather than the IN / VCC / GND some
% modules print. With VCC in the middle the pulldown, which bridges
% IN and GND, would have to cross the +5 V wire.
\draw[block, fill=white] (9.5,0.4) rectangle (13.8,3.4);
\node[font=\small\bfseries, anchor=north] at (11.65,3.3) {Relay module};
\node[note, anchor=north, align=center] at (11.65,2.45)
  {opto-isolated\\active-HIGH, N.O.\\5\,V coil};
\node[pinlbl, anchor=west] at (9.62,3.0) {IN};
\node[pinlbl, anchor=west] at (9.62,2.0) {GND};
\node[pinlbl, anchor=west] at (9.62,1.0) {VCC};
\node[pinlbl, anchor=east] at (13.68,3.0) {COM};
\node[pinlbl, anchor=east] at (13.68,2.0) {NO};
\node[pinlbl, anchor=east] at (13.68,1.0) {NC};

\draw[w, wblack]  (9.5,2.0) -- (8.4,2.0);
\draw[vw, wblack] (8.4,2.0) -- (8.4,-2.4);
\draw (8.4,-2.4) node[circ, color=wblack]{};
\draw[w, wred]    (9.5,1.0) -- (8.8,1.0);
\draw[vw, wred]   (8.8,1.0) -- (8.8,-1.6);
\draw (8.8,-1.6) node[circ, color=wred]{};

% 10k pulldown — a discrete resistor soldered across the IN and GND
% pins of the header, so it is drawn there rather than hung off a
% box edge that has no pin on it.
\draw[w, wblack] (8.9,3.0) -- (9.3,3.0);
\draw[w, wblack] (8.9,3.0) to[R] (8.9,2.0);
\draw (9.3,3.0) node[circ, color=wbrown]{};
\draw (8.9,2.0) node[circ, color=wblack]{};
\node[pinlbl, anchor=east] at (8.72,2.62) {10\,k$\Omega$};
\node[note, anchor=west, align=left, tssafe] at (9.95,-0.45)
  {The 10\,k$\Omega$ pulldown across \texttt{IN} and \texttt{GND}\\
   is \textbf{mandatory}. A floating GPIO26 — on boot,\\
   crash, or brown-out — must leave the relay OPEN.};

% ---- contacts out to the mains side ----
\draw[mainswire] (13.8,3.0) -- (16.0,3.0);
\draw[mainswire] (13.8,2.0) -- (16.0,2.0);
\draw[wire]      (13.8,1.0) -- (14.9,1.0);
\draw (16.0,3.0) node[ocirc, color=tsmains]{};
\draw (16.0,2.0) node[ocirc, color=tsmains]{};
\draw (14.9,1.0) node[ocirc]{};
\node[note, anchor=west, tsmains] at (16.2,3.0) {$\rightarrow$ thermostat NC lead};
\node[note, anchor=west, tsmains] at (16.2,2.0) {$\rightarrow$ M1 coil terminal A1};
\node[note, anchor=south] at (14.5,1.15) {leave open};

\draw[isoline] (15.5,0.5) -- (15.5,3.8);
\node[danger, anchor=south west] at (15.15,3.92) {MAINS POTENTIAL};

% ============================================================
% Colour legend
% ============================================================
\node[legendbox, anchor=north west, text width=7.9cm] at (-10.9,-4.2) {%
  \textbf{Wire colour code} — same as \texttt{WIRING.md}\\[5pt]
  \begin{tabular}{@{}p{1.7cm}@{\hspace{5pt}}p{5.4cm}@{}}
    \textcolor{wred}{\rule[0.6ex]{1.5cm}{1.3pt}}    & +5\,V rail from the PSU\\[3pt]
    \textcolor{worange}{\rule[0.6ex]{1.5cm}{1.3pt}} & 3V3, regulated on the DevKitC\\[3pt]
    \textcolor{wblack}{\rule[0.6ex]{1.5cm}{1.3pt}}  & GND / common bus\\[3pt]
    \textcolor{wyellow}{\rule[0.6ex]{1.5cm}{1.3pt}} & SPI SCK, and the K-type $(+)$ lead\\[3pt]
    \textcolor{wgreen}{\rule[0.6ex]{1.5cm}{1.3pt}}  & SPI MISO (MAX31855 \texttt{SO})\\[3pt]
    \textcolor{wblue}{\rule[0.6ex]{1.5cm}{1.3pt}}   & SPI CS\\[3pt]
    \textcolor{wpurple}{\rule[0.6ex]{1.5cm}{1.3pt}} & ADC1 input (GPIO34)\\[3pt]
    \textcolor{wbrown}{\rule[0.6ex]{1.5cm}{1.3pt}}  & relay drive (GPIO26)\\[3pt]
    \textcolor{wgrey}{\rule[0.6ex]{1.5cm}{1.3pt}}   & ack button\\[3pt]
    \textcolor{tsmains}{\rule[0.6ex]{1.5cm}{1.5pt}} & \textcolor{tsmains}{\textbf{mains potential}}\\
  \end{tabular}
};

\node[legendbox, anchor=north west, text width=8.3cm] at (-2.4,-4.2) {%
  \textbf{Reading this sheet}\\[4pt]
  A filled dot is the \emph{only} thing that means connected. Where a
  vertical wire crosses a horizontal one it breaks over the top and
  nothing is joined.\\[6pt]
  \textbf{Traps worth re-reading before power-up}\\[4pt]
  \begin{tabular}{@{}p{0.35cm}@{}p{7.4cm}@{}}
    1 & Both ESP32 GND pins land on the same rail. Jumper them.\\[3pt]
    2 & MAX31855 \texttt{VCC} is 3V3. A 5\,V feed destroys the part.\\[3pt]
    3 & The NTC must stay on GPIO34 (ADC1). ADC2 pins read zero
        while Wi-Fi is up.\\[3pt]
    4 & K-type is polarised, and \textbf{red is negative}. That is
        not a typo.\\[3pt]
    5 & The 10\,k$\Omega$ pulldown on relay \texttt{IN} is not
        optional.\\
  \end{tabular}
};

% ============================================================
% Title block
% ============================================================
\node[titleblock, anchor=north west, text width=6.6cm] at (7.2,-4.2) {%
  {\normalsize\bfseries Enclosure interior — pictorial}\\[2pt]
  {\color{tsmute}Module placement and wire colours}\\[5pt]
  {\color{tsmute}Table-saw thermal protection retrofit}\\
  {\color{tsmute}Spatial only. Pin truth: \texttt{WIRING.md}}\\
  {\color{tsmute}Topology: \texttt{esp32\_supervisor.svg}}\\
  {\color{tsmute}Source: \texttt{esp32\_pictorial.tex}}%
};

\end{circuitikz}
\end{document}
